\section{RaRa ja Digar}\label{sec:first-section}
Eesti Rahvusraamatukogu (RaRa) on avalik-õiguslik institutsioon, mis tegutseb Eesti Rahvusraamatukogu seaduse alusel ning mille põhiülesandeks on tagada kõigile Eesti elanikele juurdepääs nii digitaalsele kui füüsilisele teabele ja kultuuripärandile. Lisaks füüsilistele kogudele pakub RaRa registreeritud kasutajatele ligipääsu mitmesugustele rahvusvahelistele teadusandmebaasidele, sealhulgas väljaspool raamatukogu ruume \cite{rara}.

RaRa üks keskseid digitaalseid ressursse on digitaalarhiiv DIGAR, mis on Eesti trükiväljaannete laiapõhjaliseim digitaalne kogu \cite{rara}. Arhiiv hõlmab materjale alates esimestest Eestis ilmunud ajalehtedest ning alates 2018. aastast on digitaalsel kujul kättesaadavad kõik Eestis ilmuvad trükised. Euroopa Regionaalarengu Fondi rahastatud projekti „Kultuuripärandi digiteerimine 2018--2023" elluviimise tulemusena lisandus arhiivi üle 1,5 miljoni lehekülje tekstimaterjali, millest ligikaudu 700\,000 lehekülge ajalehti moodustab enamuse nõukogudeaegne ajakirjandus \cite{digar}.

DIGARi portaal sisaldab praegu üle 758\,858 väljaannet, 7\,441\,057 lehekülge ning 22\,797\,684 artiklit, mida täiendatakse iga päev. Portaal hõlmab 1772 ajalehte, 698 ajakirja ning 566 jätkväljaannet, kusjuures ajalehed on kättesaadavad alates 1811. aastast \cite{digar}.

\section{Digiteeritud ajalehed ja OCR}\label{sec:second-section}
Eesti ajalehtede digitaliseerimine sai alguse 1993. aastal, mil raamatukogud hakkasid haprale paberile trükitud ajalehti mikrofilmidele salvestama. Esimesed ajalehed kanti digitaalsetele piltidele 2003. aastal ning 2004. aastal avati avalik veebikeskkond DEA (Digiteeritud Eesti Ajalehed). Praeguseks on eri digiarhiivides peamiselt DEA-s, Eteras, TÜ DSpace'is ja Eesti Kirjandusmuuseumis Kivikeses kokku üle 700\,000 digiteeritud ajalehenumbri \cite[lk~75--80]{tinits2025}.

Ajalehtede digiteerimine toimub tavaliselt skaneeritud piltide tekstiks teisendamise teel, kasutades optilise märgituvastuse ehk OCR (optical character recognition) tehnoloogiat. OCR võimaldab muuta pildina esitatud teksti masinloetavaks, kuid selle kvaliteet sõltub tugevalt nii algmaterjali seisukorrast kui ka kasutatud tarkvarast. Üks kesksemaid probleeme digiteeritud ajalehtede töötlemisel ongi just OCR-i ebaühtlane kvaliteet: ajalooliste ajalehtede puhul esineb sageli valesti tuvastatud tähemärke, katkisi või poolikuid sõnu, tekstiplokkide valet järjestust ning mitme veeru tekstide soovimatut ühendamist (Laura Nemvalts, suuline kommunikatsioon, 16.09.2025).

Van Strien jt on süstemaatiliselt hinnanud OCR-kvaliteedi mõju mitmele NLP-ülesandele, Sealhulgas lausete segmenteerimisele, nimeolemite tuvastusele, sõltuvusanalüüsile, teemamudeldamisele ning keelemudelite kohandamisele ja leidnud, et OCR-vead avaldavad järjepidevat mõju kõigile testitud ülesannetele, kuigi eri ülesanded on erineva tundlikkusega \cite{vanstrien2020}.

Peeter Tinits toob välja, et DEA kogus on OCR-i kindluse hinnang lisatud vaid 62\%-le ajalehenumbritest ning ajavahemikul 1850--1950 digiteeritud materjali keskmine usaldusväärsus jääb 70--80\% juurde, samas kui 21.sajandi materjalides ulatub see 98--99\%-ni. Seega ei ole digiteeritud ajalehed ühtselt puhastatud ega ühtse kvaliteediga andmekogum, vaid pigem on tegemist ebaühtlase materjaliga, kus kvaliteet sõltub tugevalt digiteerimise ajast ja kasutatud tehnoloogiast \cite[lk~99--101]{tinits2025}.

See ebaühtlus on käesoleva lõputöö jaoks eriti oluline eeltingimus. Töö keskendub neile DIGAR-i ajalehtedele, mis on kättesaadavad vaid lehekülje kaupa, mitte artikliteks jaotatud kujul. P. Tinits märgib, et artikliteks segmenteeritud ajalehed moodustavad DEA kogust vaid osa ning enamik vanemast materjalist on salvestatud just lehekülje kaupa \cite[lk~102--103]{tinits2025}. Lehekülje põhised materjalid on reeglina rohkemate OCR-vigadega kui artikliteks jaotatud ajalehed, sest need on digiteeritud üle kümne aasta tagasi, mil kasutatud OCR-tööriistad ei vastanud praegustele standarditele. Lisaks on osa materjalist digiteeritud mikrofilmidelt, mis toob kaasa täiendava pildikvaliteedi languse ning suurendab veelgi tärktuvastuse veavõimalust (Laura Nemvalts, kirjalik kommunikatsioon, 31.03.2026).

Seega on töö andmestiku peamised väljakutsed kolmekordsed: OCR-müra tähemärgi tasandil, puuduv artiklipiiri informatsioon ning ebaühtlane kvaliteet eri digiteerimisperioodide vahel. Need eeltingimused on otseselt mõjutanud käesolevas töös valitud lähenemist, täpsemalt andmete eeltöötluse ulatust ning mudelile esitatavate sisendite kujundamist.

\section{Artiklipiiri tuvastamine digiteeritud ajalehtedes}\label{sec:third-section}
Digiteeritud ajalehtede tekstitöötlusel on üks kesksemaid väljakutseid artiklite eraldamine lehekülje tekstimassiivist. Ajaleheküljel paikneb tavaliselt mitu temaatiliselt erinevat artiklit, mis on trükitud kõrvuti veergudes ning mille piirid ei ole masinloetavas OCR-tekstis otseselt märgistatud. Ilma artiklipiiri informatsioonita on keeruline ehitada usaldusväärseid tekstianalüüsi mudeleid, kuna leheküljepõhine tekst ei moodusta temaatiliselt terviklikke üksusi \cite[lk~102]{tinits2025}.

Senised lähenemised artiklitest eraldamisele jagunevad peamiselt kahte rühma. Esimene rühm tugineb lehekülje paigutusinformatsioonile (layout analysis). Girdhar jt (2023) hindavad oma töös NewsEye artiklitest eraldamise andmestikku (NAS), mis hõlmab skaneeritud ajaleheküljed 19. ja 20.sajandist saksa, soome ja prantsuse keeles. Autorid rõhutavad, et erinev paigutus ja kirjastiilid muudavad üldistamise keeruliseks, sest mudel, mis töötab ühel ajalehel hästi, ei pruugi teisele üle kanduda \cite{girdhar2023}.

Sellele probleemile on oma lahenduse käinud välja Sun jt (2024). Nende idee seisneb LIAS-meetodi kasutamisel, mis tuvastab ajalehe eraldusjooned, analüüsib paigutusinfot ning segmenteerib teksti tekstiplokkide semantiliste seoste alusel ja lõpuks teostab artiklite eraldamise. Enamikul juhtudel prantsuse ja soome ajalehtede puhul, on autorid näidanud, et olemasolevaid lahendusi on LIAS suutnud ületada. Lehekülje visuaalse struktuuri kättesaadavus, jääb siiski mõlemal lähenemisel praktikas piiratuks \cite{sun2024}. P. Tinits (2025) märgib, et DEA kogus on artikliteks segmenteeritud vaid natukene üle poole ajalehenumbritest ning enamik vanemast materjalist on salvestatud lehekülje kaupa, ilma artiklistruktuuri informatsioonita \cite[lk~102--103]{tinits2025}.

Teise lähenemisena on uuritud suurte keelemudelite (Large Language Models, LLM) kasutamist artiklitest eraldamiseks, tuginedes ainult OCR-tekstile ilma paigutusinformatsioonita. Mauermann ja Oberbichler (2025) katsetasid seda lähenemist mitkeelsetel ajaloolistel ajalehtedel ning saavutasid mõningate mudelite ja proovide kombinatsiooniga üle 85\% täpsuse \cite{mauermann2025}. Autorid toovad siiski välja olulise piirangu: OCR-vead, eriti kohanimes ja võtmesõnades, võivad mudeli täpsust märkimisväärselt langetada. Lisaks on suurte keelemudelite põhine lähenemine arvutuslikult ressursimahukas ning mudelite käitumine raskesti ennustatav, sest erinevad mudelid reageerivad proovide sõnastusele erinevalt \cite{mauermann2025}.

Käesolevas töös on valitud hoopis muu lähenemine ehk BERT-põhine sekventsklassifikatsioon lausetasandil. Erinevalt paigutusanalüüsil põhinevatest meetoditest ei eelda see lähenemine visuaalse struktuuriinfo olemasolu, mis teeb selle sobivaks just lehekülje kaupa salvestatud materjaliga töötamiseks. Erinevalt LLM-põhisest lähenemisest ei püüa mudel mõista kogu artikli sisu tervikuna, vaid tuvastab piire lokaalselt lausetasandil: iga lause klassifitseeritakse kas artikli alguseks (silt 1) või artikli jätkuks (silt 0). See binaarne klassifitseerimisülesanne on arvutuslikult efektiivsem ning sobib hästi olukorras, kus andmestik sisaldab palju OCR-müra ja struktuurilisi ebatäpsusi, mis on käesoleva töö andmestiku peamised eeltingimused.

Mudeli aluseks on EstBERT (tartuNLP/EstBERT), mis on eesti keelel eeltreenitud BERT-arhitektuuriga mudel. BERT (Bidirectional Encoder Representations from Transformers) on keeleesituse mudel, mille eripäraks on kahesuunaline kontekstianalüüs, kus mudel arvestab iga sõna tõlgendamisel nii eelnevat kui järgnevat teksti, mis annab paremad tulemused semantiliselt keerulistes klassifitseerimisülesannetes võrreldes ühepoolsete mudelitega \cite{devlin2019}. Selline arhitektuur võimaldab eeltreenitud mudelit kohandada paljudele erinevate NLP-ülesannete jaoks minimaalse ülesandespetsiifilise muutusega \cite{wolf2020}. Käesolevas töös kasutatakse seda omadust artiklipiiri tuvastamiseks ehk peenhäälestuse käigus kohandatakse eeltreenitud mudel segmenteeritud ajalehtedest koostatud treeningandmestiku põhjal, kus artiklistruktuur on teada.
